\subsection{Написание кеша для метаданных}

\subsubsection{Задача}
\begin{itemize}
  \item Реализован read model кеш, наполняемый при старте сервиса из слоя персистенции.
  \item Кеш подключается как отдельная библиотека с минимальным набором зависимостей. 
  \item Инвалидация кеша осуществляется по событиям через event listener без перезапуска сервиса. 
  \item При обращении к данным из кеша количество обращений к базе данных происходит только на заполнение кеша.
  \item Кеш корректно обновляется при получении события изменения метаданных. 
\end{itemize}

\subsubsection{Аналитика задачи}

Данная задача была описана аналитиком после общих идей по введению кеширования в подсистемы. 
Система корректно работает только при обычной нагрузке, но при повышении нагрзуки в 10 раз,
подсистемы излишне много делают запросов в однокластерную персистенцию. 
Две подсистемы для формульных преобразований и для трансформаций файлов так же активно ходят в персистенцию, но я заметил,
что от полученных данных они используют только 30 процентов информации. 
Для данных двух подсистем можно реализовать некоторый общий инструмент,
который позволит на старте приложения вытягивать данные из персистенции, 
сохранять их в оперативную память, что повысит локальность данных и скорость их получения.
Данные архитектор попросил предоставить в упрощённом варианте, так как полноценные документы из документоориентированной базы будут излишне занимать место.
Первичная нагрузка ожидается в 50.000 элементов по 30 полей в каждом. Поля есть как служебные, так и обычные.
Read модель кеша означает то, что его принято только читать, 
а при изменениях он будет инвалидироваться и перезаполняться, 
так как изменения этой части данных будут изменяться не часто -- константы.
Поскольку кеш кодово будет одинаковым для обоих сервисов, то решено сделать его ввиде подключаемой библиотеки, с набором интерфейсов для кастомной реализации, если потребуется.

\subsubsection{Реализация решения}

Приступив к реализации, я для начала решил определить базовую структуру библиотеки.
Поскольку она должна подключаться к двум независимым сервисам,
я оформил её как отдельный Maven-модуль (см. Приложение Г)
с зависимостями на $Spring Context$, $Jackson$ и клиент $Nuxeo$,
остальное библиотека может взять после подключения к конкретному сервису,
так как меня попросили её сделать максмаильно легковесной. 
Данная библиотекта будет иметь версию без указания $SNAPSHOT$.

\subsubsection{Модель данных}
После создания проекта спроектировал упрощённую модель элемента справочника —
класс $DictionaryItem$ (Приложение Г). Так как архитектор обозначил требование хранить только 30\% полей документа,
я выделил набор служебных полей -- $\_id$,
$\_label$, $\_value$, $\_parent$, $\_is\_default$,
а оставшиеся поля разместил в словаре с помощью аннотаций
Jackson $JsonAnyGetter$ $JsonAnySetter$, что позволит мне не 
привязывать модель кеша к конкретной схеме документов базы, в которой есть около 100 излишних полей.

\subsubsection{Слой получения данных}
Для получения данных из персистенции я ввёл интерфейс
$RdmData$ $Provider$, а его реализацию $RdmDataProviderRestImpl$ делегирует HTTP-клиенту
$DataProviderRdmClient$ (Приложение Г).
Клиент обращается к эндпоинту с Basic-аутентификацией
и таймаутом 30 минут, так как первичная выгрузка элементов
может занимать продолжительное время, а со временем ожидается только увеличение размера.
Так же был написан NuxeoRdmService (Приложение Г) который позволяет получать уникальный сгенерированный номер документа для дальнейшего запроса иерарихии документов с главного узла
по его ссылке из персистенции.

\subsubsection{Парсинг иерархии справочников}
Данные из персистенции поступают в виде вложенной JSON-структуры, 
так как между персистенцией и библиотекой стоит ещё один сервис, 
который приводит данные в транспортировочный формат.

Я определил интерфейс $RdmParserService$
и реализовал $RdmParser$ $ServiceStringImpl$ (Приложение Г).
Алгоритм является рекурсивным, но далее по возможности и времени можно его разложить в линейный,
он обходит JSON-дерево, ищет ключи с суффиксом который задан в $RdmClientConfig$,
извлекает по ним имена справочников и разворачивает вложенные массивы
в плоскую структуру, так как иерархичность для данных сервисов не играет роли.

\subsubsection{Кеш и инициализация}
Главным компонентом стал сервис $RdmLibService$ (Приложение Г), он практически как оркестратор.
Для хранения данных кеша я выбрал примитив $ConcurrentHashMap$ со значениями
типа $CopyOnWriteArrayList$ для обеспечения безопасности
при параллельных чтениях без доп блокировок. После первичного тестирования я обнаружил, 
что не стоит допускать пользователя до кеша, пока он не заполнился, так как возникают ошибки системы из-за пустых данных.
Чтобы потребитель не мог обратиться к кешу до завершения его наполнения,
я ввёл барьер на базе $CountDownLatch$ с таймаутом 10 минут:

\begin{lstlisting}[language=Java]
private final CountDownLatch cacheInitLatch = new CountDownLatch(1);

public List<DictionaryItem> getDictByDocRef(String docRef) {
    cacheInitLatch.await(10, TimeUnit.MINUTES);
    ...
}
\end{lstlisting}

\subsubsection{Асинхронная инициализация и повторные попытки}
Заполнение кеша запускается по кастомному событию $RdmRefresh$ $Event$
(Приложение Г), которое сервис-потребитель публикует при старте приложения 
(я сделал специально некоторое делегирование этой функции на сторону пользователя библиотеки, так как у раных сервисов понятие полного старта может отличаться).
Слушатель $RefreshEventListener$ (Приложение Г) помечен аннотацией
для асинхронного запуска в отдельном потоке, 
поэтому не блокирует поток запуска контекста.
При неудаче загрузки слушатель выжидает некоторое заданное вреся
и повторяет попытку -- это покрывает случаи временной
недоступности персистенции при старте кластера, 
недоступность персистенции связана с заполнением кешей документоориентированной прослойки.

\subsubsection{Конфигурация}
Для подключения к Nuxeo я оформил отдельный конфиг
$NuxeoRdm$ $Config$ (Приложение Г),
читающий параметры подключения из $application$ $.properties$
сервиса-потребителя.
Все интерфейсы
вынесены в публичное API библиотеки,
чтобы при необходимости сервис мог
подменить реализацию. 
Я сделал это для того, чтобы не менять код библиотеки, если потребуется какая-то узконаправленная функция.

\subsubsection{Тестирование}
Покрытие я реализовал в RdmLibServiceTest (Приложение Г),
проверив три сценария: полноценный разбор вложенной структуры,
разбор минимального справочника и обработку null-значений в полях.
Тесты используют JSON тестовые файлы (\texttt{single\_dictionary.json, null\_dictionary.json} -- Приложение Г).

\subsubsection{Интеграция решения}

После публикации библиотеки в корпоративный Maven-репозиторий                
я подключил её к обоим сервисам одинаковым образом:                     
в \texttt{pom.xml} ввиде зависимости: 
                                                                              
\begin{lstlisting}[language=XML]                                             
<dependency>                                                                 
    <groupId>ru.rabis.utils.rdm</groupId>                                    
    <artifactId>rdm-manager</artifactId>                                     
</dependency>               
\end{lstlisting}                                                             
                
Библиотека поставляет собственные бины, поэтому пришлось добавить                  
в аннотации \texttt{ComponentScan} главного класса каждого приложения       
пакет библиотеки, чтобы контекст обнаружил и зарегистрировал их автоматически.                 
                                                                              
Для запуска наполнения кеша в каждом сервисе реализован                      
\texttt{Command LineRunner}, что и является отправной точкой валидного старта.
После завершения старта контекста они публикуют событие                      
\texttt{RdmRefreshEvent}, тем самым передавая управление                     
асинхронному слушателю библиотеки, который начинает заполнение кеша.       
                                                                              
Есть некоторый набор классов, которые являются некоторыми функциями, 
которые вызываются из шаблона трансформаций,     
каждая из них обращается к кешу                   
через \texttt{rdmLibService.getDictByDocRef(rdmName)},                       
минуя персистенцию. 

При получении сигнала об изменении                       
справочных данных сервис публикует повторный \texttt{RdmRefreshEvent},       
что запускает инвалидацию и перезагрузку кеша. 
Ранее эндпоинт был уже интегрирован в сервис, но инвалидации касались немного другого кеша.                        
                                                                                                               
Интеграция свелась к трём шагам для каждого сервиса:          
добавить зависимость, расширить \texttt{@ComponentScan} и                    
опубликовать \texttt{RdmRefreshEvent} при старте, инвалидации.                            
Остальное --- наполнение кеша, барьер готовности, повторные попытки ---      
взяла на себя библиотека.

\subsubsection{Заполенения кеша}
После старта приложения \texttt{CommandLineRunner} публикует событие
\texttt{RdmRefreshEvent("Init")}.
Логи при успешной инициализации:

\begin{lstlisting}[language={}]
- Refresh event: Init
- rdm-library:1.0.1
- Vocabularies upload size: 49857
- Init success
\end{lstlisting}

Строка \texttt{Vocabularies upload size: 49857} подтверждает, что все
справочники получены и помещены в \texttt{ConcurrentHashMap}.
Далее защёлка снимется и у сервиса будет доступ к справочникам.

\subsubsection{Чтения из кеша}
В метод \texttt{getDictByDocRef} я не стал по умолчанию вписывать логгирование,
так как он вызывается очень часто. Для верификации поведения я временно добавил
строку логирования:

\begin{lstlisting}[language=Java]
public List<DictionaryItem> getDictByDocRef(String docRef) {
    awaitCacheInit();
    if (docRef != null && !docRef.isBlank()) {
        CopyOnWriteArrayList<DictionaryItem> list = cache.get(docRef);
        log.info("Cache lookup '{}': {} items", docRef,
                  list != null ? list.size() : 0);
        return list != null ? new ArrayList<>(list) : null;
    }
    log.error("docRef blank");
    return null;
}
\end{lstlisting}

При чтении кеша в лог выводится:

\begin{lstlisting}[language={}]
Cache lookup 'Absceng': 52 items
Cache lookup 'PriceRoute': 18 items
Cache lookup 'LaborPrice': 34 items
\end{lstlisting}

\subsubsection{Инвалидации кеша}
При изменении справочных данных сервис публикует повторный
\texttt{RdmRefreshEvent("Evict")}. Слушатель срабатывает асинхронно:
в \texttt{refreshCache()} вызывается \texttt{cache.clear()}. Логи инвалидации:

\begin{lstlisting}[language={}]
- Refresh event: Evict
- rdm-library:1.0.1
- Vocabularies upload size: 49857
- Init success
\end{lstlisting}

Последовательность идентична первичной инициализации, как раз поэтому мы видим одинаковое кол-во загруженных элементов.
\subsubsection{Снижения обращений в базу}
HTTP-запрос к персистенции производится только внутри
\texttt{get Hierarhy()}, который вызывается
исключительно из \texttt{refreshCache()}.
Все последующие обращения к справочникам обслуживались кешем.
Это означает, что cache hit rate составил 100\%.

\subsubsection{Вывод}
Реализовал read model кеш, который наполняется исключительно при старте сервиса из слоя персистенции.
Логически кеш реализован ввиде библиотеки с минимальным набором зависимостей, что выделяет легковестность его зависимости. 
Инвалидация кеша сделаны через событийную модель через event listener без перезапуска сервиса. 
При обращении к данным из кеша количество обращений к базе данных происходит только на заполнение кеша.